\section{Problem Formulation}
\label{sec:problem}

We model one agent step as a branching step served by a prefill instance $\Pin$ and a decode instance $\Din$ connected by a KV connector.
\Cref{tab:notation} in the appendix summarizes the notation.

\begin{definition}[Branching step]
\label{def:step}
A parent $u$ holds a committed trunk $x_u$ with $|x_u|{=}L$.
A fan-out is a set of candidate continuations.
We call the candidate-specific suffix $\rho_i$ a \emph{residual}, with $|\rho_i|{=}\ell_i$; candidate $i$ is the sequence $x_u\Vert\rho_i$.
Index $0$ denotes the highest-ranked candidate produced by the branching controller (e.g., the greedy continuation).
Commit selects one continuation $\star$ and leaves the spine $x_\star=x_u\Vert\rho_\star$, which becomes the trunk of the next step.
\end{definition}

With $b$ bytes of KV per token, three occupancies characterize one step:
\begin{align}
M_{\mathrm{clone}} &= b\bigl(kL+\textstyle\sum_i\ell_i\bigr), \label{eq:mem}\\
M_{\mathrm{CoW}} &= b\bigl(L+\textstyle\sum_i\ell_i\bigr), \nonumber\\
M_{\mathrm{spine}} &= b\bigl(L+\ell_\star\bigr). \nonumber
\end{align}
$M_{\mathrm{clone}}$ arises only when every request materializes its own trunk; $M_{\mathrm{CoW}}$ holds one shared trunk and all residuals; $M_{\mathrm{spine}}$ is the occupancy after uncommitted residuals are freed.

\subsection{Residency on Two Instances}
\label{sec:residency}

Paged engines index full blocks of $P$ tokens by a chained key $h_j=\mathrm{hash}(h_{j-1},\,\text{block}_j,\,\text{extra})$~\cite{vllm,sglang}, so a block is matched only together with its entire parent prefix.
For an instance $\mathcal{X}\in\{\Pin,\Din\}$ we define the resident prefix of the trunk and of residual $i$ as
\begin{align}
c_u^{\mathcal{X}} &= \max\{jP\le L:\ x_u[0{:}jP)\sqsubseteq\mathcal{X}\}, \label{eq:mi}\\
m_i^{\mathcal{X}} &= \max\{jP\le \ell_i:\ x_u\Vert\rho_i[0{:}jP)\sqsubseteq\mathcal{X}\}, \nonumber
\end{align}
where $z\sqsubseteq\mathcal{X}$ means that every page of $z$ is physically resident on $\mathcal{X}$,
and $m_i^{\mathcal{X}}{=}0$ whenever $c_u^{\mathcal{X}}{<}L$, because chained keys make a residual page unreachable without its full trunk.
For exposition we take $L$ to be a multiple of $P$; otherwise the partial trunk page is charged to each residual.
Three distinctions matter.
(i)~\emph{Parent-aware match:} $m_i^{\mathcal{X}}$ is defined on $x_u\Vert\rho_i$, not on $\rho_i$ alone.
(ii)~\emph{Index presence versus residency:} an index entry whose block has been evicted does not count; $m_i^{\mathcal{X}}$ counts physically resident pages only.
(iii)~\emph{Instance locality:} pages resident on $\Pin$ are not visible on $\Din$, so $m_i^{\Pin}$ and $m_i^{\Din}$ generally differ.

\subsection{Cache Regimes}
\label{sec:regimes}

Prefix caching alone does not determine whether a fan-out pays for the trunk once or $k$ times.
Let
\begin{equation}
\label{eq:lbar}
\bar L_{\mathcal{X}}=L-c_u^{\mathcal{X}},\qquad \mathcal{X}\in\{\Pin,\Din\},
\end{equation}
be the trunk tokens missing on instance $\mathcal{X}$.
We define $\nu_u^{\Pin}$ as the number of distinct materializations of $\bar L_{\Pin}$ on the prefill instance and $\nu_u^{\Din}$ as the number of distinct connector transfers of $\bar L_{\Din}$.
Thus
\begin{equation}
\label{eq:amp}
W^{\mathrm{trunk}}=\nu_u^{\Pin}\bar L_{\Pin},\qquad
X^{\mathrm{trunk}}=\nu_u^{\Din}\bar L_{\Din},
\end{equation}
with $\nu_u^{\mathcal{X}}{=}0$ if $\bar L_{\mathcal{X}}{=}0$, and $1\le\nu_u^{\mathcal{X}}\le k$ otherwise.
We additionally let $\sigma_u\in\{0,1\}$ denote a publication barrier: $\sigma_u{=}1$ delays sibling admission until one reusable trunk copy is published.
A cache regime is therefore characterized by
\begin{equation}
\label{eq:regime}
\mathcal{R}_u=(\nu_u^{\Pin},\nu_u^{\Din},\sigma_u).
\end{equation}
The prefill-side regimes are special cases:
\begin{equation}
\label{eq:rlist}
\begin{aligned}
R_1\ \text{(cold, concurrent):}\ &
  c_u^{\Pin}{=}0,\ \sigma_u{=}0,\\
&\quad 1\le\nu_u^{\Pin}\le k,\\
R_2\ \text{(cold, serialized):}\ &
  c_u^{\Pin}{=}0,\ \sigma_u{=}1,\ \nu_u^{\Pin}{=}1,\\
R_3\ \text{(warm):}\ &
  c_u^{\Pin}{=}L,\ W^{\mathrm{trunk}}{=}0,\\
R_4\ \text{(explicit fork):}\ &
  \sigma_u{=}0,\ \nu_u^{\Pin}{=}1.
\end{aligned}
\end{equation}
Hence $R_1$ gives $L\le W^{\mathrm{trunk}}\le kL$, $R_2$ gives $W^{\mathrm{trunk}}{=}L$ but introduces a serialization barrier, and $R_4$ gives $W^{\mathrm{trunk}}{=}L-c_u^{\Pin}$ without serialization.
Connector behavior is characterized independently by $\nu_u^{\Din}$: request-local transfer may have $\nu_u^{\Din}{=}k$, whereas shared-trunk aliasing gives $\nu_u^{\Din}{=}1$.

For $\bar L_{\mathcal{X}}>0$, $\nu_u^{\mathcal{X}}$ is the trunk-amplification factor.
Relative to independent fan-out, sharing saves $(k-\nu_u^{\Pin})\bar L_{\Pin}$ prefill tokens and $(k-\nu_u^{\Din})\bar L_{\Din}$ connector tokens.
All GPU admission experiments in \S\ref{sec:eval} use $R_3$ on $\Pin$, so their prefill savings are attributable to residual admission rather than trunk sharing.
After commit, the spine is the logical trunk of the next step; its KV residency on $\Pin$ and $\Din$ is restored according to the underlying KV synchronization policy.

\subsection{Objective}
\label{sec:objective}

Each residual receives an action $a_i$ from a set $\Act$, inducing computed prefill $W_i$, connector tokens $X_i$, and decode tokens $d_i$; we write $W{=}\sum_iW_i$, $X{=}\sum_iX_i$, and $D^{+}{=}\sum_id_i$.
We formulate the design objective as
\begin{equation}
\label{eq:objective}
\begin{aligned}
\min_{\pi}\;\; & \E\bigl[\Pref(W)+\mu X+T_{\mathrm{dec}}(D^{+})\bigr]\\
\text{s.t.}\;\; & \Prob[\text{committed answer correct}]\ge Q_{\min},\\
& M_{\mathrm{peak}}\le M_{\mathrm{pool}},\qquad a_0\neq\Reject,
\end{aligned}
\end{equation}
where $\mu$ is the per-token transfer time.
The defining constraint is informational: a decision taken before prefill may depend only on $(c_u^{\Pin},c_u^{\Din},m_i^{\Pin},m_i^{\Din})$ and the residual text, whereas a decision taken on $\Din$ may additionally read generated tokens.
Decode-time pruners act only at the second point; \sys adds the first and coordinates the two.
\sys instantiates this objective with a lightweight threshold-based policy.
